% chktex-file 1
% chktex-file 3
% chktex-file 8
% chktex-file 36
\documentclass[11pt]{article}

\usepackage[margin=1in]{geometry}
\usepackage{amsmath,amssymb,amsthm}
\usepackage{bm}
\usepackage{xspace}
\usepackage{booktabs}
\usepackage{enumitem}
\usepackage{listings}
\usepackage{xcolor}
\usepackage[round,authoryear]{natbib}
\usepackage{hyperref}
\usepackage[capitalize]{cleveref}
\hypersetup{colorlinks=true, linkcolor=blue, urlcolor=blue, citecolor=blue}

\newcommand{\bw}{\mathbf{w}}
\newcommand{\bs}{\mathbf{s}}
\newcommand{\bY}{\mathbf{Y}}
\newcommand{\bI}{\mathbf{I}}
\newcommand{\bZero}{\boldsymbol{0}}
\newcommand{\skewt}{skew-$t$\xspace}
\newcommand{\Skewt}{Skew-$t$\xspace}

\newcommand{\eref}[1]{(\ref{#1})}
\newcommand{\sref}[1]{Section~\ref{#1}}

\theoremstyle{plain}
\newtheorem{proposition}{Proposition}[section]

\theoremstyle{definition}
\newtheorem{definition}[proposition]{Definition}

\theoremstyle{remark}
\newtheorem{remark}[proposition]{Remark}
\newtheorem*{intuition}{Intuition}

\lstset{
    basicstyle=\ttfamily\small,
    keywordstyle=\color{blue},
    commentstyle=\color{gray}\itshape,
    stringstyle=\color{teal},
    showstringspaces=false,
    breaklines=true,
    frame=single,
    framesep=4pt,
    language=R
}

\title{AR(2) MCMC for Beginners:\\
       A Walkthrough of Issue~1 and Issue~2}
\author{Yi-Xuan Xie}
\date{2026-05-15}

\begin{document}
\maketitle

\begin{abstract}
This note is a beginner-friendly companion to
\texttt{tex/ar2\_fix\_proposal/} and \texttt{tex/where\_we\_are/}.
Those earlier notes assume the reader is comfortable with the Metropolis--Hastings
algorithm, detailed balance, and the Markov factorisation of an AR(2) joint
prior. Here we build the same machinery from the ground up and then return
to the two issues that motivated the audit, explaining \emph{why} each one
is real and \emph{how} the patch removes it, using explicit mathematical
notation throughout. We treat the third issue (the $\phi$ Hastings
adjustment) elsewhere; that one is more subtle and is not the focus here.
\end{abstract}

\tableofcontents

\section{Why this note}\label{sec:why}

The AR(2) extension of the \citet{morris_space-time_2017} space-time \skewt
model places a stationary AR(2) prior on three transformed parameters at
each spatial knot. Updating these parameters inside an MCMC requires care
because the AR(2) prior links each time point to the next two, and the
update step has to reflect all of that linkage. Two issues arise:

\begin{description}
  \item[Issue~1.] The initial pair $(Y_1, Y_2)$ of a stationary AR(2)
        series is not independent --- the two values share a covariance
        $\gamma_1$ implied by the recursion. The simulator in the code
        respects this; the MCMC density at $t = 2$ \emph{previously} did
        not, and has now been patched to use the stationary form.
  \item[Issue~2.] When updating a single time point $Y_t$, the Metropolis--Hastings
        acceptance ratio must include \emph{three} prior factors of the
        joint distribution. The original code included only two. The third,
        missing factor (the ``lag-2 term'') has been added at all four
        call sites.
\end{description}
Both patches are now in production and guarded by regression tests in
\texttt{code/R/ar2/tests/}; this note explains \emph{why} each patch was
needed and \emph{how} the math works out.

This note builds up just enough background --- stationarity, Yule--Walker,
the MH algorithm, detailed balance --- to make both issues unambiguous in
math. We then state each patch as an explicit formula.

\section{Background: the AR(2) prior}\label{sec:bg}

Throughout this section, $\{Y_t\}_{t=1}^{n_t}$ is a generic scalar series
that stands in for any of the transformed parameters used in the model
($\bw^{*}_{tk}$ componentwise, $z^{*}_{tk}$, or $\sigma^{2*}_{tk}$, at a
fixed spatial knot~$k$).

\subsection{The recursion}\label{sec:recursion}

The AR(2) model says that each $Y_t$, for $t \geq 3$, is determined by
the two previous values plus a fresh Gaussian shock:
\begin{equation}\label{eq:recursion}
  Y_t \;=\; \phi_1\, Y_{t-1} \;+\; \phi_2\, Y_{t-2} \;+\; \epsilon_t,
  \qquad
  \epsilon_t \stackrel{iid}{\sim} N(0, \sigma_\epsilon^2),
  \qquad t \geq 3.
\end{equation}
Two parameters $(\phi_1, \phi_2)$ describe the linear dependence on the
past. One parameter $\sigma_\epsilon^2$ describes the noise level. Two
initial values $(Y_1, Y_2)$ kick the recursion off. That is the entire
model.

\subsection{What ``stationary'' means}\label{sec:stationary}

A time series is (weakly) stationary if three things hold for every $t$:
\begin{align*}
  \text{(i) } E[Y_t] &= \text{constant in } t, \\
  \text{(ii) } \mathrm{Var}(Y_t) &= \text{constant in } t, \\
  \text{(iii) } \mathrm{Cov}(Y_t, Y_{t+h}) &= \gamma_h,
        \text{ a function of the lag } h \text{ only}.
\end{align*}
For \eref{eq:recursion} to define a stationary process, the parameters
$(\phi_1, \phi_2)$ must lie inside the so-called \emph{stationarity
triangle}:
\begin{equation}\label{eq:stab}
  |\phi_2| < 1, \qquad \phi_1 + \phi_2 < 1, \qquad \phi_2 - \phi_1 < 1.
\end{equation}

\begin{intuition}
Stationarity is what lets us talk about ``the marginal distribution of
$Y_t$'' without specifying which $t$ we mean. In our application this is
crucial: the paper relies on each $Y_t$ being marginally $N(0, 1)$ so
that the inverse copula transform produces the right \skewt distribution
at every time point.
\end{intuition}

\subsection{Yule--Walker: deriving the moments}\label{sec:moments}

We want the stationary process to have unit variance, $\gamma_0
:= \mathrm{Var}(Y_t) = 1$. Combined with stationarity, this constraint
forces specific values on the autocovariances $\gamma_1, \gamma_2, \ldots$
and on the noise variance $\sigma_\epsilon^2$. We derive them now.

Multiply both sides of \eref{eq:recursion} by $Y_{t-h}$ (for some lag
$h \geq 1$) and take expectations:
\begin{equation}\label{eq:yw}
  \gamma_h \;=\; \phi_1\, \gamma_{h-1} \;+\; \phi_2\, \gamma_{h-2}.
\end{equation}
This is the \emph{Yule--Walker} equation. It does not involve
$\sigma_\epsilon^2$ because $\epsilon_t$ is independent of $Y_{t-h}$ for
$h \geq 1$.

\paragraph{Lag $h = 1$.}
Plug $h = 1$ into \eref{eq:yw}, and use the symmetry
$\gamma_{-1} = \gamma_1$:
\[
  \gamma_1
  \;=\; \phi_1\, \gamma_0 + \phi_2\, \gamma_{-1}
  \;=\; \phi_1 + \phi_2\, \gamma_1
  \;\Longrightarrow\;
  \boxed{\;\gamma_1 = \dfrac{\phi_1}{1 - \phi_2}.\;}
\]

\paragraph{Lag $h = 2$.}
Plug $h = 2$ into \eref{eq:yw}:
\[
  \boxed{\;\gamma_2 \;=\; \phi_1\, \gamma_1 + \phi_2.\;}
\]

\paragraph{Innovation variance.}
Multiply \eref{eq:recursion} by $Y_t$ (not $Y_{t-h}$) and take
expectations. Now $E[Y_t \epsilon_t] = E[\epsilon_t^2] = \sigma_\epsilon^2$
because the same-time shock is in $Y_t$. We get
\[
  \gamma_0 \;=\; \phi_1\, \gamma_1 + \phi_2\, \gamma_2 + \sigma_\epsilon^2
  \;\Longrightarrow\;
  \boxed{\;\sigma_\epsilon^2 \;=\; 1 - \phi_1\, \gamma_1 - \phi_2\, \gamma_2.\;}
\]

\begin{intuition}
Yule--Walker is the bridge between the two ways of describing an AR(2):
the \emph{recursion} (in terms of $\phi_1, \phi_2, \sigma_\epsilon^2$)
and the \emph{covariance structure} (in terms of $\gamma_0, \gamma_1,
\gamma_2$). Once we fix $\gamma_0 = 1$, the rest is determined by
$(\phi_1, \phi_2)$ alone.
\end{intuition}

\subsection{Innovation variance: AR(1) for comparison}\label{sec:ar1-vs-ar2}

Setting $\phi_2 = 0$ collapses AR(2) to AR(1):
\[
  \gamma_1 \to \phi_1, \qquad \gamma_2 \to \phi_1^2, \qquad
  \sigma_\epsilon^2 \to 1 - \phi_1^2.
\]
This matches the AR(1) formulation in the paper exactly. It also explains
why the AR(1) version of the model has none of the subtleties we are
about to discuss in Issues~1 and~2: when there is no $\phi_2$, $Y_t$ only
links to one neighbour ($Y_{t-1}$) instead of two.

\section{Issue 1: the stationary joint of \texorpdfstring{$(Y_1, Y_2)$}{(Y1, Y2)}}\label{sec:issue1}

\subsection{The recursion needs two seeds}\label{sec:two-seeds}

\eref{eq:recursion} only kicks in at $t = 3$ -- to compute $Y_3$, we need
$Y_2$ and $Y_1$. So we have to specify, separately, the joint
distribution of $(Y_1, Y_2)$. Marginal statements like ``$Y_1 \sim N(0,1)$''
and ``$Y_2 \sim N(0,1)$'' do not pin this down --- they leave the covariance
$\mathrm{Cov}(Y_1, Y_2)$ undetermined.

\subsection{What stationarity forces on \texorpdfstring{$(Y_1, Y_2)$}{(Y1, Y2)}}\label{sec:stationary-joint}

If we want the full series $Y_1, \ldots, Y_{n_t}$ to be a stationary AR(2),
the initial pair has to follow the stationary joint distribution implied
by $\gamma_0 = 1$ and $\gamma_1 = \phi_1/(1-\phi_2)$:
\begin{equation}\label{eq:joint}
  \boxed{\;
  \begin{pmatrix} Y_1 \\ Y_2 \end{pmatrix}
  \;\sim\;
  N_2\!\left(
    \begin{pmatrix} 0 \\ 0 \end{pmatrix},\;\;
    \begin{pmatrix} 1 & \gamma_1 \\ \gamma_1 & 1 \end{pmatrix}
  \right).
  \;}
\end{equation}
Marginally, both $Y_1$ and $Y_2$ are $N(0, 1)$; jointly they correlate at
strength $\gamma_1$. If you initialised them as independent $N(0, 1)$ and
then ran the recursion forward with the Yule--Walker $\sigma_\epsilon^2$,
the resulting $Y_3, Y_4, \ldots$ would \emph{not} have unit variance --- the
chain would only stabilise to a stationary regime after a transient,
violating the marginal property at every $t$ during the transient.

\subsection{The conditional density implied by \eref{eq:joint}}\label{sec:cond}

When the MCMC updates $Y_2$ alone, it needs the conditional $Y_2 \mid Y_1
= y_1$. Apply the standard bivariate Gaussian conditioning formula to
\eref{eq:joint}:
\begin{equation}\label{eq:p2-stationary}
  \boxed{\;
  Y_2 \mid Y_1 = y_1 \;\sim\; N\bigl(\gamma_1\, y_1,\; 1 - \gamma_1^2\bigr).
  \;}
\end{equation}
The conditional mean is $\gamma_1 y_1$, not $\phi_1 y_1$; the conditional
variance is $1 - \gamma_1^2$, not $\sigma_\epsilon^2$. We will call this
``the stationary $t=2$ density'' below.

\subsection{What the code uses instead}\label{sec:code-t2}

The function \texttt{get\_ar2\_conditional\_params} in
\texttt{auxfunctions\_ar2.R} returns, at $t = 2$:
\begin{equation}\label{eq:p2-code}
  Y_2 \mid Y_1 = y_1 \;\sim\; N\bigl(\phi_1\, y_1,\; \sigma_\epsilon^2\bigr).
\end{equation}
This is the form one would get by pretending the recursion already
applies at $t = 2$ (with a fictitious lag-2 value $Y_0 = 0$).

\subsection{Why they differ unless \texorpdfstring{$\phi_2 = 0$}{phi2 = 0}}\label{sec:why-differ}

Compare \eref{eq:p2-stationary} and \eref{eq:p2-code}:
\begin{center}
\begin{tabular}{lll}
\toprule
                  & Mean of $Y_2 \mid Y_1$ & Variance of $Y_2 \mid Y_1$ \\
\midrule
Stationary form   & $\gamma_1 y_1 = \dfrac{\phi_1}{1-\phi_2}\, y_1$
                  & $1 - \gamma_1^2 = 1 - \left(\dfrac{\phi_1}{1-\phi_2}\right)^2$ \\
Code's form       & $\phi_1\, y_1$
                  & $\sigma_\epsilon^2 = 1 - \phi_1 \gamma_1 - \phi_2 \gamma_2$ \\
\bottomrule
\end{tabular}
\end{center}
The two coincide when $\phi_2 = 0$ (since then $\gamma_1 \to \phi_1$ and
$\sigma_\epsilon^2 \to 1 - \phi_1^2 = 1 - \gamma_1^2$). For genuinely
AR(2) parameters they differ.

\subsection{Bug or modelling choice?}\label{sec:bug-vs-choice}

It depends on which target model you adopt:

\begin{itemize}
  \item \emph{Target~A (stationary AR(2)):} for every $t$, $Y_t$ is
        marginally $N(0,1)$. The $t = 2$ conditional is
        \eref{eq:p2-stationary}. Under this target, the code is wrong at
        $t = 2$.
  \item \emph{Target~B (independent initial values, then AR(2) recursion):}
        $Y_1, Y_2$ are independent $N(0,1)$, then $Y_t$ for $t \geq 3$
        follows \eref{eq:recursion}. Under this target, the $t = 2$
        conditional is degenerate (no lag, just the marginal $N(0,1)$),
        and \eref{eq:p2-code} is wrong too.
\end{itemize}
Neither target is dictated by the paper, which only says $Y_1, Y_2 \sim
N(0,1)$ marginally. The code's choice \eref{eq:p2-code} doesn't fit
either target cleanly; it is a halfway-house. The simulator
\texttt{simulate\_ar2\_standard} does follow Target~A. So at $t = 2$
the simulator and the MCMC use different conditionals.

The impact under the previous code was one mis-specified factor out of
$n_t - 1$, i.e.\ a relative bias of order $1/n_t$. We have since adopted
Target~A: the $t = 2$ branch of \texttt{get\_ar2\_conditional\_params}
now returns the stationary form \eref{eq:p2-stationary}. With that
one-line change the simulator and the MCMC density agree on the joint
of $(Y_1, Y_2)$, and the unit marginal variance property
$\mathrm{Var}(Y_t) = 1$ holds at every $t$. Regression coverage is
provided by \texttt{tests/test\_fix1\_t2\_conditional.R}, which checks
the returned $(\mathrm{mean}, \mathrm{sd})$ against
$(\gamma_1 y_1, \sqrt{1 - \gamma_1^2})$ to machine precision at six
$(\phi_1, \phi_2)$ configurations.

\section{The Metropolis--Hastings update of \texorpdfstring{$Y_t$}{Yt}}\label{sec:mh}

We now build the framework needed for Issue~2.

\subsection{What problem does MCMC solve?}\label{sec:problem}

Bayesian inference says: combine the likelihood $L(\bY \mid \text{parameters})$
with a prior $p(\text{parameters})$ to get the posterior
$p(\text{parameters} \mid \bY) \propto L \cdot p$. For complex models
this posterior usually does not have a closed-form density we can sample
from, so we build an MCMC sampler instead. The sampler produces a Markov
chain whose stationary distribution \emph{is} the posterior. Run the
chain long enough, and the samples it produces are samples from the
posterior.

\subsection{Gibbs sampling and full conditionals}\label{sec:gibbs}

A common strategy --- Gibbs sampling --- is to update one parameter (or one
block of parameters) at a time, drawing from its \emph{full conditional}:
the conditional distribution of that parameter given all the other
parameters and the data. For some parameters, the full conditional has
a closed-form density (e.g.\ a conjugate Gaussian or Gamma). For others
it does not, and we have to use Metropolis--Hastings as a substitute
within the Gibbs cycle.

In our model, the time-varying parameters $z^{*}_{tk}$ (and analogous
$\bw^{*}_{tk}, \sigma^{2*}_{tk}$) do not have closed-form full
conditionals because the AR(2) prior is convolved with the spatial
likelihood. So we Metropolis them.

\subsection{The Metropolis--Hastings algorithm}\label{sec:mh-algo}

Let $\pi(Y)$ denote the target density (here: the full conditional of
$Y_t$). The MH algorithm is:

\begin{enumerate}[label=(\arabic*)]
  \item Start the chain at some value $Y$.
  \item Propose a candidate $Y^{(c)}$ from a \emph{proposal density}
        $q(Y^{(c)} \mid Y)$. A common choice is a Gaussian random walk,
        $Y^{(c)} \sim N(Y, s^2)$.
  \item Compute the acceptance probability
        \begin{equation}\label{eq:alpha}
          \alpha(Y, Y^{(c)})
          \;=\; \min\!\left(1,\;\;
            \frac{\pi(Y^{(c)})\, q(Y \mid Y^{(c)})}
                 {\pi(Y)\, q(Y^{(c)} \mid Y)}
          \right).
        \end{equation}
  \item With probability $\alpha$, set $Y \leftarrow Y^{(c)}$; otherwise,
        $Y$ stays.
  \item Repeat from step (2).
\end{enumerate}
After a burn-in, the values of $Y$ along the chain are dependent samples
from $\pi$.

\subsubsection*{Symmetric proposal: simplification}

When $q$ is symmetric, $q(Y^{(c)} \mid Y) = q(Y \mid Y^{(c)})$ (Gaussian
random walks are symmetric), the proposal ratio cancels and the
acceptance probability reduces to
\begin{equation}\label{eq:alpha-sym}
  \alpha(Y, Y^{(c)})
  \;=\; \min\!\left(1,\; \frac{\pi(Y^{(c)})}{\pi(Y)}\right).
\end{equation}
This is the form the code uses for the $Y_t$ updates.

\subsection{Detailed balance: why MH works}\label{sec:db}

The mathematical guarantee behind MH is \emph{detailed balance}: a chain
with transition kernel $P(\cdot \to \cdot)$ has $\pi$ as its stationary
distribution if, for all states $y$ and $y^{(c)}$,
\begin{equation}\label{eq:db}
  \pi(y)\, P(y \to y^{(c)}) \;=\; \pi(y^{(c)})\, P(y^{(c)} \to y).
\end{equation}
With MH, $P(y \to y^{(c)}) = q(y^{(c)} \mid y)\, \alpha(y, y^{(c)})$, and
the acceptance rule in \eref{eq:alpha} is the unique rule that makes
\eref{eq:db} hold. So if we plug in the \emph{right} $\pi$, the chain
samples from $\pi$.

\subsection{The right \texorpdfstring{$\pi$}{pi} for the \texorpdfstring{$Y_t$}{Yt} update}\label{sec:right-pi}

For Gibbs-within-MH, the right $\pi$ is the full conditional density of
$Y_t$ given all other parameters and the data:
\[
  \pi(Y_t)
  \;:=\; p(Y_t \mid Y_{-t}, \bY, \text{other params})
  \;\propto\; L(\bY \mid Y_t) \cdot p(Y_t \mid Y_{-t}).
\]
The right side has two pieces: the data likelihood $L$ and the
conditional prior of $Y_t$ given the rest of the time series.

To get the second piece, write down the joint prior of the whole series
$(Y_1, \ldots, Y_{n_t})$. Under the AR(2) Markov structure, this joint
factorises as
\begin{equation}\label{eq:joint-prior}
  p(Y_1, \ldots, Y_{n_t})
  \;=\; p(Y_1)\, p(Y_2 \mid Y_1)\,
        \prod_{s=3}^{n_t} p(Y_s \mid Y_{s-1}, Y_{s-2}).
\end{equation}
The conditional prior of a single $Y_t$ given the rest is proportional
to whatever factors of \eref{eq:joint-prior} depend on $Y_t$. Other
factors don't depend on $Y_t$ and cancel as constants in the ratio.

\paragraph{Which factors involve $Y_t$?}
Looking at \eref{eq:joint-prior}, the value $Y_t$ shows up in three
factors (for $3 \leq t \leq n_t - 2$):
\begin{itemize}
  \item $p(Y_t \mid Y_{t-1}, Y_{t-2})$ --- $Y_t$ as the ``self''
        conditioned variable.
  \item $p(Y_{t+1} \mid Y_t, Y_{t-1})$ --- $Y_t$ as the lag-1 input.
  \item $p(Y_{t+2} \mid Y_{t+1}, Y_t)$ --- $Y_t$ as the lag-2 input.
\end{itemize}
All other factors (for $s \notin \{t, t+1, t+2\}$) do not contain $Y_t$
and cancel. So the full conditional is
\begin{equation}\label{eq:full-cond}
  \boxed{\;
  \pi(Y_t \mid Y_{-t}, \bY)
  \;\propto\;
  L(\bY \mid Y_t)\,\cdot\,
  \underbrace{p(Y_t \mid Y_{t-1}, Y_{t-2})}_{\text{self}}
  \,\cdot\,
  \underbrace{p(Y_{t+1} \mid Y_t, Y_{t-1})}_{\text{lag-1}}
  \,\cdot\,
  \underbrace{p(Y_{t+2} \mid Y_{t+1}, Y_t)}_{\text{lag-2}}.
  \;}
\end{equation}

\begin{remark}
For $t = n_t - 1$ the lag-2 factor is absent (there is no $Y_{n_t+1}$).
For $t = n_t$ both the lag-1 and lag-2 factors are absent. Edge cases
matter when implementing the patch.
\end{remark}

\section{Issue 2: the missing lag-2 factor in the acceptance ratio}\label{sec:issue2}

We are now ready to state Issue~2 precisely.

\subsection{What the buggy code computes}\label{sec:buggy-R}

The original \texttt{updateZTS\_AR2} (and analogous functions for the
other parameters) computes the acceptance ratio as a log-quantity
\[
  \log R_{\mathrm{buggy}}
  \;=\; \log\frac{L(\bY \mid Y_t^{(c)})}{L(\bY \mid Y_t)}
        \;+\;\Delta_{\text{self}}
        \;+\;\Delta_{\text{lag-1}},
\]
where
\begin{align*}
  \Delta_{\text{self}}  &= \log p(Y_t^{(c)} \mid Y_{t-1}, Y_{t-2})
                          - \log p(Y_t \mid Y_{t-1}, Y_{t-2}), \\
  \Delta_{\text{lag-1}} &= \log p(Y_{t+1} \mid Y_t^{(c)}, Y_{t-1})
                          - \log p(Y_{t+1} \mid Y_t,     Y_{t-1}).
\end{align*}
The lag-2 factor $p(Y_{t+2} \mid Y_{t+1}, Y_t)$ is missing.

\subsection{Why this is wrong: target shift}\label{sec:target-shift}

Plug $R_{\mathrm{buggy}}$ into the MH acceptance rule \eref{eq:alpha-sym}.
By detailed balance \eref{eq:db}, the resulting chain has as its
stationary distribution not the full conditional \eref{eq:full-cond} but
the modified density
\begin{equation}\label{eq:wrong-target}
  \widetilde\pi(Y_t)
  \;\propto\;
  L(\bY \mid Y_t)\,\cdot\,
  p(Y_t \mid Y_{t-1}, Y_{t-2})\,\cdot\,
  p(Y_{t+1} \mid Y_t, Y_{t-1}),
\end{equation}
in which the lag-2 prior factor is silently removed. $\widetilde\pi$ is
\emph{not} the AR(2) full conditional. The chain produces samples whose
stationary covariance structure does not match $\gamma_2 = \phi_1 \gamma_1
+ \phi_2$ --- the lag-2 autocorrelation drifts away from the Yule--Walker
value.

\subsection{The patch in math}\label{sec:patch-math}

The fix is to add the missing factor:
\begin{equation}\label{eq:patch}
  \boxed{\;
    \log R_{\mathrm{patched}}
    \;=\; \log R_{\mathrm{buggy}}
        \;+\;\Delta_{\text{lag-2}},
  \;}
\end{equation}
where
\begin{equation}\label{eq:delta-lag2}
  \Delta_{\text{lag-2}}
  \;=\;
  \log p(Y_{t+2} \mid Y_{t+1}, Y_t^{(c)})
  - \log p(Y_{t+2} \mid Y_{t+1}, Y_t).
\end{equation}

\subsubsection{Each Gaussian factor explicitly}

The lag-2 factor is Gaussian:
\[
  p(Y_{t+2} \mid Y_{t+1}, Y_t)
  \;=\; \frac{1}{\sqrt{2\pi\sigma_\epsilon^2}}
        \exp\!\Bigl(-\tfrac{1}{2\sigma_\epsilon^2}\,
                    (Y_{t+2} - \phi_1 Y_{t+1} - \phi_2 Y_t)^2\Bigr).
\]
The log-difference \eref{eq:delta-lag2} simplifies (the
$\log(2\pi\sigma_\epsilon^2)$ piece cancels because it does not depend on
$Y_t$):
\begin{align}
  \Delta_{\text{lag-2}}
  &= \frac{1}{2\sigma_\epsilon^2}
     \Bigl[(Y_{t+2} - \phi_1 Y_{t+1} - \phi_2 Y_t)^2
         - (Y_{t+2} - \phi_1 Y_{t+1} - \phi_2 Y_t^{(c)})^2\Bigr] \notag\\
  &= \frac{\phi_2\bigl(Y_t^{(c)} - Y_t\bigr)
           \bigl[\,2(Y_{t+2} - \phi_1 Y_{t+1}) - \phi_2(Y_t + Y_t^{(c)})\,\bigr]}
          {2\sigma_\epsilon^2}.
  \label{eq:delta-explicit}
\end{align}

\begin{intuition}
The factor $\Delta_{\text{lag-2}}$ acts like a ``future anchor'' for
$Y_t$. When the proposal $Y_t^{(c)}$ moves in a direction that makes the
predicted value $\phi_1 Y_{t+1} + \phi_2 Y_t^{(c)}$ closer to the
observed $Y_{t+2}$, the bracketed term in \eref{eq:delta-explicit}
has the right sign to boost acceptance. Without this anchor, the chain
treats $Y_t$ as if $Y_{t+2}$ were not influenced by it --- a falsehood
under the AR(2) recursion --- and the lag-2 autocorrelation of the chain
ends up wrong. Note that $\Delta_{\text{lag-2}} \propto \phi_2$: when
$\phi_2 = 0$, the anchor vanishes (AR(2) collapses to AR(1)).
\end{intuition}

\subsection{The patch in code}\label{sec:patch-code}

In R, the patch is a single block inserted right after the existing
lag-1 correction:
\begin{lstlisting}
# Fix 2: Y_t enters the t+2 conditional density as lag-2
if ((t + 2) <= nt) {
    z.star.next2 <- z.star[k, t + 2]
    R <- R +
        ar2_conditional_log_density(
            z.star.next2, t + 2,
            z.star[k, t + 1], can.z.star[k],     # propose Y_t^(c) as lag-2
            phi1, phi2, moments) -
        ar2_conditional_log_density(
            z.star.next2, t + 2,
            z.star[k, t + 1], z.star[k, t],      # current Y_t as lag-2
            phi1, phi2, moments)
}
\end{lstlisting}
The same insertion appears at four sites in
\texttt{code/R/ar2/update\_params\_ar2.R}: two branches of
\texttt{updateTauTS\_AR2}, one in \texttt{updateZTS\_AR2}, and one in
\texttt{updateKnotsTS\_AR2} (the last with a \texttt{sum(\ldots)} over
the $K \times 2$ knots array).

\subsection{What the test shows}\label{sec:test}

A prior-only chain (no data likelihood, just the AR(2) prior) was run
with truth $(\phi_1, \phi_2) = (0.7, -0.4)$, giving theoretical
$\gamma_1 = 0.500$ and $\gamma_2 = -0.050$. The empirical
autocorrelations after burn-in:
\begin{center}
\begin{tabular}{lrrrr}
\toprule
chain                   & mean      & var       & lag-1 ACF & lag-2 ACF \\
\midrule
target                  & $0$       & $1.000$   & $+0.500$  & $-0.050$  \\
buggy   (no $t+2$ term) & $+0.009$  & $1.142$   & $+0.608$  & $+0.222$  \\
patched (Fix~2 applied) & $+0.005$  & $0.998$   & $+0.498$  & $-0.054$  \\
\bottomrule
\end{tabular}
\end{center}
The buggy chain inflates the marginal variance by $14\%$ and gets the
lag-2 autocorrelation \emph{wrong in sign}; the patched chain matches
the Yule--Walker target to four decimal places. This empirically confirms
\eref{eq:wrong-target}: removing the lag-2 factor shifts the target away
from the AR(2) full conditional.

\section{Summary}\label{sec:summary}

Two issues, both about the joint structure of the AR(2) prior:

\begin{description}
  \item[Issue~1.] The stationary AR(2) implies
        $(Y_1, Y_2) \sim N_2(\bZero, [\![1, \gamma_1; \gamma_1, 1]\!])$,
        with the conditional
        $Y_2 \mid Y_1 \sim N(\gamma_1 Y_1, 1 - \gamma_1^2)$ \eref{eq:p2-stationary}.
        The code's MCMC previously used $N(\phi_1 Y_1, \sigma_\epsilon^2)$
        \eref{eq:p2-code} at $t = 2$ instead, which is the form of the
        AR(2) recursion extrapolated backward with a fictitious $Y_0 = 0$;
        the two differ unless $\phi_2 = 0$. The $t = 2$ branch of
        \texttt{get\_ar2\_conditional\_params} has been replaced by
        \eref{eq:p2-stationary}, restoring the simulator-MCMC agreement
        on the joint of $(Y_1, Y_2)$ and preserving
        $\mathrm{Var}(Y_t) = 1$ at every $t$. Status: \textbf{patched,
        tested} (\texttt{tests/test\_fix1\_t2\_conditional.R}).

  \item[Issue~2.] The MH acceptance ratio for $Y_t$ has to include
        \emph{three} factors of the AR(2) joint prior --- self
        ($p(Y_t \mid Y_{t-1}, Y_{t-2})$), lag-1
        ($p(Y_{t+1} \mid Y_t, Y_{t-1})$), and lag-2
        ($p(Y_{t+2} \mid Y_{t+1}, Y_t)$). The original code computes only
        the first two. The chain therefore samples from a target
        $\widetilde\pi$ \eref{eq:wrong-target} that has the wrong lag-2
        autocorrelation. The patch \eref{eq:patch}--\eref{eq:delta-lag2}
        restores the missing factor at four call sites, and a prior-only
        regression test confirms the corrected chain matches
        Yule--Walker. Status: \textbf{patched, tested}.
\end{description}

Issue~1 is about the \emph{initial conditions} of the AR(2) recursion;
Issue~2 is about the \emph{Metropolis--Hastings ratio} during ongoing
updates. They are independent in spirit; one can hold without the other.
The AR(1) case has neither issue because (a) the AR(1) recursion
automatically lands in its stationary joint when started from
$Y_1 \sim N(0,1)$, and (b) $Y_t$ only ever influences $Y_{t+1}$ (not
$Y_{t+2}$), so the AR(1) acceptance ratio with two factors is complete.

\bibliographystyle{plainnat}
\bibliography{../references/references}

\end{document}
